Aus $\frac{1}{2}m\overline{v^2} = \frac{3}{2}kT$ (FoTaBe S.~171) folgt $v_\text{rms} = \sqrt{3kT/m}$. Erdgas ist ein Gemisch (vor allem Methan), deshalb verwendet man seine Dichte. Mit der Teilchenzahldichte $N/V$ ist $\rho = \frac{N}{V}\,m$, und die ideale Gasgleichung $p = \frac{N}{V}\,kT$ ergibt
\begin{align}
 \frac{kT}{m} &= \frac{p}{\rho}
\end{align}
Mit dem Normaldruck $p_0 = \ncpo$ und der Dichte von Erdgas bei Normalbedingungen $\rho = \rEg$ (typischer Wert, FoTaBe S.~187) wird
\begin{align}
 v_\text{rms} &= \vEgF = \sqrt{\frac{3\cdot\ncpo}{\rEg}} = \vEg \approx \result{\vEgP{0}{2}}
\end{align}
Für reines Methan ($\rho = \rMe$) ergäbe sich $\vMeP{0}{2}$.

Bemerkung: Die mittlere Schnelligkeit $\bar v = \sqrt{8kT/(\pi m)}$ (FoTaBe S.~171) ist etwas kleiner, für Erdgas $\bar v = \vbEgF = \vbEgP{0}{2}$.
